\chapter*{Jak się uczyć z tej książki}
\addcontentsline{toc}{chapter}{Jak się uczyć z tej książki}
\markboth{Jak się uczyć}{Jak się uczyć}

\section*{Rytm dnia: tryb 100 (domyślny)}

\begin{center}
\begin{tabular}{@{}rl@{}}
\toprule
\textsf{\textbf{Czas}} & \textsf{\textbf{Co robisz}}\\
\midrule
15 min & \emph{Anki} --- powtórki (zawsze najpierw), potem maks.\ 15--20 nowych kart\\
25 min & nowy materiał z jednostki (słownictwo, dialog, reguła)\\
20 min & ćwiczenia: pisz odpowiedzi \emph{ręcznie lub w edytorze}, nie w głowie\\
20 min & słuchanie lub czytanie autentyczne (DW, wiadomości, podcast)\\
20 min & mówienie na głos albo pisanie: rozmowa z AI lub korepetytorem\\
\bottomrule
\end{tabular}
\end{center}

\section*{Rytm tygodnia}

\begin{center}
\begin{tabular}{@{}lp{0.72\linewidth}@{}}
\toprule
\textsf{\textbf{Dzień}} & \textsf{\textbf{Fokus}}\\
\midrule
Poniedziałek & słownictwo i dialog jednostki (czytaj na głos, trzy razy)\\
Wtorek & gramatyka, część pierwsza + ćwiczenia 1--4\\
Środa & gramatyka, część druga + ćwiczenia 5--koniec\\
Czwartek & słuchanie i czytanie: te same struktury w autentycznym tekście\\
Piątek & mówienie i pisanie: rozmowa na temat jednostki, krótki tekst\\
Sobota & zadanie weekendowe (produkcja), korekta z AI lub korepetytorem\\
Niedziela & lekko: tylko Anki + godzina słuchania \emph{bez} napisów po polsku\\
\bottomrule
\end{tabular}
\end{center}

\section*{Zasady, które decydują o wyniku}

\begin{zasada}[Reguła 1: produkcja]
Co tydzień \emph{mów} i \emph{pisz} to, czego się uczysz. Zrozumienie tekstu to nie umiejętność
mówienia. Jeśli czytasz dialog, a potem go nie odtwarzasz z pamięci, wyćwiczyłeś rozpoznawanie.
\end{zasada}

\begin{zasada}[Reguła 2: zdania, nie reguły]
Gramatykę uczysz się w całych zdaniach. Zapamiętaj \de{Ich lerne seit einem Jahr Deutsch}, a nie
tabelkę przypadków. Tabela potem wynika ze zdań, nie odwrotnie.
\end{zasada}

\begin{zasada}[Reguła 3: nie przechodzisz dalej bez 80\%]
Ćwiczenia sprawdzasz z kluczem w Dodatku~A. Poniżej 80\% poprawnych odpowiedzi powtarzasz
część gramatyczną, zanim przejdziesz do kolejnej jednostki. To dlatego są tygodnie buforowe.
\end{zasada}

\begin{zasada}[Reguła 4: rodzajnik razem z rzeczownikiem]
Zawsze \de{der Tisch}, nigdy samo \de{Tisch}. Fiszki w \code{anki/} już tak działają.
\end{zasada}

\section*{Jak używać Anki}

W katalogu \code{anki/} jest po jednym pliku na poziom (\code{german-a1.tsv} itd.). Import:
\emph{File $\rightarrow$ Import}, plik jest już opisany nagłówkami (typ notatki \emph{Basic and
reversed card}, talia \code{Deutsch::A1}, tagi z numerem jednostki). Importuj jednostki \emph{w
miarę postępu}, nie wszystko naraz: 15--20 nowych kart dziennie to sufit, nie cel.

\section*{Jak używać asystenta AI do rozmów}

Poproś o rozmowę na poziomie, na którym jesteś, i wymagaj: (1) poprawiania każdego błędu
z wyjaśnieniem, (2) odpowiedzi wyłącznie po niemiecku na poziomie A1/A2/B1/B2, (3) pytań
zwrotnych. Przykład polecenia:

\begin{uwaga}[Polecenie startowe]
\de{Du bist mein Deutschlehrer auf Niveau A2. Sprich nur Deutsch, kurze Sätze. Stelle mir Fragen
zum Thema dieser Woche. Korrigiere jeden Fehler und erkläre ihn auf Polnisch in einem Satz.}
\end{uwaga}

\section*{Pułapki w samej nauce}

\begin{itemize}
\item \textbf{Zmiana materiałów co tydzień.} Wybierz jeden kurs główny (ta książka + DW) i trzymaj się go.
\item \textbf{Iluzja płynności.} Jeśli zawsze zerkasz na tłumaczenie, nie uczysz się. Zasłoń kolumnę.
\item \textbf{Perfekcjonizm w mówieniu.} Błędy w rozmowie są paliwem. Cel to komunikacja.
\item \textbf{Dialekt za wcześnie.} Dialekt dopiero od B1, jako rozumienie (Dodatek~D).
\end{itemize}
